% =====================================================================
% Benchmarking and metrics for poorly structured LLM-generated code
% Kai Yi Ng -- Machine Learning Research Group, University of Oxford
%
% Compile: pdflatex main -> bibtex main -> pdflatex main -> pdflatex main
% (Overleaf does this automatically on Recompile.)
% =====================================================================
\documentclass[11pt,a4paper]{article}

\usepackage[utf8]{inputenc}
\usepackage[T1]{fontenc}
\usepackage[margin=1in]{geometry}
\usepackage{booktabs}
\usepackage{graphicx}
\usepackage{amsmath}
\usepackage{listings}
\usepackage{xcolor}
\usepackage[hidelinks]{hyperref}
\usepackage{url}

\lstset{
  basicstyle=\ttfamily\footnotesize,
  breaklines=true,
  frame=single,
  framerule=0.3pt,
  rulecolor=\color{gray!50},
  showstringspaces=false,
  language=Python,
}

\title{Which Measures Can See a Code Smell?\\
       \large A Validation Protocol for Assessing the Structural Quality
       of LLM-Generated Code}

\author{Kai Yi Ng\\
        Machine Learning Research Group\\
        University of Oxford\\
        \texttt{kaiyi.ng@some.ox.ac.uk}}

\date{\today}

\begin{document}
\maketitle

% ---------------------------------------------------------------------
\begin{abstract}
Benchmarks for code-generating language models measure whether the generated code
runs. They do not measure whether it is well built. Work that addresses the second
question reaches for the structural metrics software engineering already has ---
cyclomatic complexity, the maintainability index, Halstead measures, similarity
scores --- and reports their values on model output. Those metrics were designed
to describe human-written code, and their sensitivity to any particular structural
defect has not been established. Reporting a metric is not the same as showing it
can see what it is being used to detect.

This paper supplies the missing step. We build a benchmark of 1{,}300 Python
functions in which twelve structural defects are introduced one at a time into
otherwise unchanged code, each confirmed by a static analyser and each verified by
execution to preserve behaviour, and pair it with 4{,}300 real defective functions
drawn from human-written repositories against a shared clean baseline. Measuring
each metric on both sources answers two different questions --- whether changing
only the defect moves the metric, and whether defective and clean code differ on it
in the wild --- and detection requires yes to both.

Applying the protocol to eighteen measures, we find that no single measure detects
all twelve defects while every defect is detected by at least one, which is the
empirical case for a panel rather than a composite quality score. We find that a
substantial part of what looks like detection on real code is co-occurrence with
size: metrics that provably cannot respond to a defect still separate defective
from clean code, because defective functions are simply larger. We report a clean
negative result for language-model perplexity, which does not distinguish
structurally poor code from clean code on either source. Finally, applying the
validated panel to the output of three models, we show that scoring raw model
output conflates verbosity with structural quality: one model appeared two and a
half times more defective than another almost entirely because it volunteered test
cases the prompt did not request.
\end{abstract}

% ---------------------------------------------------------------------
\section{Introduction}

Language models now write a substantial share of new code, and the benchmarks used
to judge them ask almost exclusively whether that code runs. HumanEval and MBPP
score a generated function by executing it against unit tests; a model is better
if a larger fraction of its solutions pass. This is a reasonable thing to measure
and a poor proxy for the thing practitioners actually worry about. Code that passes
its tests can still be code no one wants to maintain --- a function that does seven
things, a conditional with six clauses, a constant with no name and no explanation.
Correctness and structural quality are different properties, and only one of them
is routinely benchmarked.

The natural response is to reach for the metrics software engineering already has.
Cyclomatic complexity, the Halstead measures, and the maintainability index have
decades of use behind them; similarity measures such as BLEU and CodeBLEU compare
generated code against a reference. Recent work does exactly this, computing these
metrics over model output and drawing conclusions about how well models write. The
step that is skipped is the one that licenses those conclusions. These metrics were
designed to \emph{describe} human-written code in aggregate, not to \emph{detect}
particular structural defects, and their sensitivity to any given defect is
assumed rather than demonstrated. A metric that moves when defective and clean code
are compared has not thereby been shown to respond to the defect --- only to
something that differs between the two groups.

That distinction is not academic, and it cannot be settled by observation alone.
Defective code found in the wild differs from clean code in many ways at once. It
is written by different people, in different projects, for different purposes --- and,
most importantly, it is bigger. Every structural metric responds to size. So a
metric can separate the two groups convincingly while being formally incapable of
responding to the defect that defines them. Distinguishing the two requires an
intervention: change only the defect, hold everything else fixed, and see whether
the metric moves. Conversely, an intervention alone is not enough either. Code with
a defect introduced deliberately is not code that acquired one naturally, and a
metric validated only on synthetic edits may be responding to the habits of
whatever produced them.

We therefore validate on two sources at once. A controlled benchmark supplies the
intervention: 1{,}300 Python functions in which each of twelve structural defects is
introduced individually into otherwise unchanged code, with the defect confirmed by
a static analyser and behaviour preservation confirmed by executing the function's
own tests. A real-world layer supplies external validity: roughly 4{,}300 defective
functions drawn from human-written repositories, compared against a common baseline
of 1{,}000 clean functions mined from the same population. A measure is credited with
detecting a defect only when it responds on both --- when the controlled experiment
shows the defect causes the movement, and the real data shows the movement is
present in code people actually wrote.

The result is a three-way verdict for every measure and defect: the measure
\emph{detects} the defect, it merely \emph{co-occurs} with it, or it is
\emph{blind} to it. Most of what a purely observational study would report as
detection falls into the middle category. We then apply the validated panel to the
output of three code-generating models, which surfaces a measurement problem that
only becomes visible with more than one model: apparent structural quality is
partly a measure of how much a model volunteers beyond what was asked.

\subsection{Contributions}

\begin{enumerate}
  \item \textbf{A validation protocol} that separates detection from correlation by
        pairing a controlled injection benchmark with real labelled code, and a
        single statistic --- detection strength --- reported on both so the two are
        directly comparable.
  \item \textbf{A labelled benchmark} of 1{,}300 Python functions covering twelve
        structural defects, every example confirmed by a static analyser and
        verified by execution to preserve behaviour, released openly.
  \item \textbf{A validated verdict for eighteen measures}, showing that no single
        measure covers all twelve defects, that every defect is covered by at
        least one, and that structural and reference-based measures are
        complementary rather than redundant.
  \item \textbf{Two cautionary results.} Language-model perplexity does not
        distinguish structurally poor code from clean code, so naturalness is not
        a proxy for structural quality; and smell rates computed on raw model
        output partly measure verbosity, which distorts multi-model comparison.
  \item \textbf{An applied evaluation} of three code models against the validated
        panel, and an open evaluation tool that applies the same panel to
        arbitrary Python code.
\end{enumerate}

\subsection{Research questions}

\begin{description}
  \item[RQ1] Which measures are sensitive to which structural defects, and does
             that sensitivity survive contact with real, human-written code?
  \item[RQ2] What is the structural profile of LLM-generated code under the
             validated panel, and does that profile hold across models?
  \item[RQ3] Do prompts written to induce a structural defect actually induce it,
             and does the number of tokens a model produces predict the structural
             quality of what it produces?
\end{description}

% ---------------------------------------------------------------------
\section{Background and related work}
% TO WRITE. Four strands, and the gap this paper occupies:
%
%   (a) Code smells: Fowler's catalogue, and detector-based identification
%       (Pylint, Ruff, jscpd). Position detectors as an operational definition,
%       not ground truth -- this is what licenses using them to label the data.
%   (b) Structural metrics: cyclomatic, Halstead, maintainability index, cognitive
%       complexity. Note they were designed as descriptive aggregates over
%       codebases, never validated as per-defect detectors. This is the gap.
%   (c) Evaluating LLM-generated code: HumanEval/MBPP and pass@1; similarity
%       measures (BLEU, CodeBLEU, CodeBERTScore); the correctness monoculture.
%   (d) Structural quality of LLM output: the emerging strand. Velasco et al.
%       measure propensity with a logit probe on already-smelly code, and state
%       as an open question whether propensity predicts the quality of freely
%       generated code -- which is what RQ2 answers.
%
% The gap sentence to land: prior work reports metric values on model output;
% none first establishes which metrics are sensitive to which defects.

% ---------------------------------------------------------------------
\section{The benchmark}
% TO WRITE.
% 3.1 The twelve defects and why these twelve (common in LLM output AND
%     mechanically confirmable -- both criteria are load-bearing).
% 3.2 Sources of clean code: MBPP (993), HumanEval (178), CodeSearchNet (129).
%     CodeSearchNet is needed because short tasks cannot host the larger defects.
% 3.3 Injection: one defect per example, worked example (Table + listing).
% 3.4 Two-stage verification: detector confirms the shape, execution confirms
%     behaviour. All twelve at 100% behaviour-preserving. Include the
%     inefficient-copy bug that only execution caught -- it makes the case for
%     the second stage better than any argument.
% 3.5 The real-code layer: the reused datasets are all-positive, so the clean
%     baseline had to be mined (1,000 methods from 253,762 unlabelled). Per-defect
%     counts table.

% ---------------------------------------------------------------------
\section{The measure panel}
% TO WRITE. Five families and why the panel excludes the detectors themselves
% (they produced the labels; evaluating them on their own labels is circular).
% Table: measure, family, reference-free or not, direction.
% The reference-free/reference-based split is the key structural fact: it
% determines which measures can be validated on real code at all.

% ---------------------------------------------------------------------
\section{Validation protocol}
% TO WRITE.
% 5.1 Detection strength: (mean_defective - mean_clean) / pooled SD, oriented so
%     positive = worse, capped at +/-5. Interpretation bands.
% 5.2 Why unpaired rather than paired -- the paired form saturates on
%     constant-shift injections (dead_code read 5.0 for one added line) and is
%     not computable on real code, which has no twins. Unpaired makes the two
%     sources directly comparable, which is the whole point.
% 5.3 The three verdicts: Strong (both sources agree), Co-occurrence (real only
%     -- the intervention shows the measure cannot see the defect), Blind.

% ---------------------------------------------------------------------
\section{RQ1: which measures detect which defects}
% TO WRITE. The main results table (injected d vs real d per defect x measure),
% the trust table, and the four disagreement patterns with worked examples:
%   - both high            -> genuine detection (deep nesting / cyclomatic)
%   - injected 0, real high -> co-occurrence (inefficient loop / cyclomatic:
%                              the rewrite provably adds no branches)
%   - injected high, real lower -> the injection was unrealistically uniform
%   - both ~0              -> blind
% Figures: detection-strength heatmap; injected-vs-real scatter.

% ---------------------------------------------------------------------
\section{RQ2: the structural profile of generated code}
% TO WRITE. Three models on 664 benchmark tasks. pass@1, defect rates, structure
% against the canonical solutions, cost. Two findings to lead with:
%   - generated code is simpler and far more heavily commented than the canonical
%     solutions, on every complexity measure, for every model;
%   - only four or five of the twelve defects appear at all, because single short
%     functions cannot host the rest. Frame as a scoped limitation of the task
%     set, and the motivation for RQ3.
% Then the verbosity confound: the as-emitted vs definitions-only split, the
% mbpp_17 example, and why this distorts any multi-model comparison that scores
% raw output.

% ---------------------------------------------------------------------
\section{RQ3: prompt-induced defects and token usage}
% TO WRITE. 476 prompts written to induce a named defect, across three models.
%   - induction rate: does a prompt engineered to provoke a defect provoke it?
%     Measurable on the six defects shared between the two taxonomies.
%   - does output length predict structural quality, controlling for the prompt's
%     stated complexity?
% NOTE for the write-up: in the prompt set, expected token depth is a perfect
% relabelling of complexity (basic/intermediate/advanced -> low/medium/high, 1:1).
% It therefore cannot serve as an independent variable; measured output tokens can.

% ---------------------------------------------------------------------
\section{Threats to validity}
% TO WRITE.
% Construct: detectors are the operational definition of a defect, not ground
%   truth; the panel excludes them to avoid circularity, but the labels remain
%   tool-derived.
% Internal: injection may not reproduce how defects arise naturally -- which is
%   precisely why the real-code layer exists; residual size confounding is
%   reported rather than corrected, and marked as co-occurrence.
% External: Python only; small instruct-tuned models plus one frontier model;
%   the similarity family is validated on injected data alone, because real code
%   has no clean reference.
% Conclusion: single greedy generation per task, so no sampling variance is
%   estimated.

% ---------------------------------------------------------------------
\section{Conclusion}
% TO WRITE. The one-sentence version: before asking what a metric says about
% LLM-generated code, establish what the metric can see -- and for most metrics
% and most defects, the answer is less than it appears.

% ---------------------------------------------------------------------
\section*{Data and code availability}
The benchmark, the evaluation tool, and every figure in this paper are available
at \url{https://github.com/kaiyi03/Code-Smells}. An interactive version of the
evaluation tool, which applies the same panel to arbitrary Python code, is at
\url{https://code-smell-dashboard.onrender.com}.

% ---------------------------------------------------------------------
\bibliographystyle{plain}
\bibliography{references}

\end{document}
